\setcounter{chapter}{7}
\chapter{Conclusion and Future Scope}
\label{ch:conclusion}

\section{Summary of Contributions}
\label{sec:summary}

This project set out to address the fragmentation that characterises modern technical recruitment by building a platform that consolidates applicant tracking, candidate evaluation, and offer management into a single web application. SmartHire achieves this through the integration of six core modules: automated ATS resume scoring, timed MCQ assessments, a browser-based coding IDE, AI-powered video interview evaluation, a recruiter Kanban pipeline, and automated PDF offer letter generation.

The platform was implemented using Next.js~15, TypeScript, Supabase PostgreSQL, the Monaco editor, and Google's Gemini API. A multi-schema database design enforces data isolation across recruiting organisations, while application-level query scoping and Supabase Row-Level Security policies provide defence-in-depth against cross-tenant data leakage.

Empirical evaluation demonstrated that automated evaluation stages reduce per-candidate processing time from over two hours of manual effort to under ten seconds of machine computation. Scoring for structured stages (ATS, MCQ, coding) is fully deterministic, eliminating inter-rater variability. Stage-specific rejection tracking---a feature absent from most commercial ATS platforms---provides candidates with transparent, actionable feedback on their application status.

\section{Lessons Learned}
\label{sec:lessons}

Several lessons emerged during the development process:

\begin{enumerate}[leftmargin=1.5em]
  \item \textbf{Prompt engineering is iterative.} The quality of AI-generated ATS scores and interview evaluations improved significantly through iterative refinement of the prompts sent to the Gemini API. Initial prompts produced inconsistent output formats; adding explicit JSON schema specifications and example outputs resolved this.

  \item \textbf{Database schema design decisions have long-term consequences.} The choice to organise tables into domain-specific schemas rather than a flat \texttt{public} schema added initial complexity but paid dividends in query clarity, access control granularity, and developer onboarding.

  \item \textbf{Client-side singletons prevent subtle bugs.} The Supabase browser client singleton pattern resolved authentication state inconsistencies that were difficult to diagnose because they only manifested during rapid page navigation.

  \item \textbf{Timer-based features require careful implementation.} The MCQ countdown timer initially drifted due to reliance on \texttt{setInterval}. Computing remaining time from wall-clock differences rather than incrementally decrementing a counter produced accurate results.
\end{enumerate}

\section{Future Scope}
\label{sec:future}

Several enhancements are planned or under consideration for future development:

\subsection{Multi-Language Code Execution}

The current coding assessment module supports JavaScript and Python. Extending support to Java, C++, Go, and Rust would require containerised execution environments (e.g., Docker-based sandboxes) to handle the diverse compilation and runtime requirements of each language.

\subsection{Advanced Proctoring}

Integrating browser-level proctoring capabilities---webcam monitoring, tab-switch detection, and clipboard restriction---would increase the integrity of remote assessments, particularly for high-stakes evaluations.

\subsection{Calendar and Email Integration}

Native integration with Google Calendar and Outlook for interview scheduling, and with email services for automated notifications, would reduce the manual coordination currently required for interview logistics.

\subsection{Mobile Applications}

A React Native mobile application would allow candidates to browse jobs, receive notifications, and accept offers from their smartphones, improving engagement and response times.

\subsection{Advanced Analytics and Reporting}

Expanding the analytics dashboard with hiring funnel conversion rates, time-to-hire tracking, source-of-hire attribution, and diversity metrics would provide recruiters with deeper insights into their hiring processes.

\subsection{Bias Detection and Mitigation}

Incorporating fairness auditing tools that analyse whether ATS scores or AI interview evaluations exhibit systematic bias against demographic groups would be an important step toward responsible AI deployment in recruitment.

\subsection{Enterprise Features}

For enterprise adoption, the platform would benefit from single sign-on (SSO) integration, role-based access control with customisable permissions, audit logging with tamper-evident storage, and compliance reporting for data protection regulations such as GDPR.

\section{Conclusion}
\label{sec:final-conclusion}

SmartHire demonstrates that it is feasible to build a unified, end-to-end recruitment platform that automates the mechanical stages of technical hiring while preserving human authority over final decisions. The platform reduces administrative overhead, produces consistent evaluation scores, and provides candidates with the transparency they deserve. While several areas remain for future development---particularly around scale testing, multi-language support, and bias auditing---the current implementation establishes a solid foundation for a production-grade recruitment automation system.
